\respcap{cap:fluxo}
\resp{fl-cortes} Com $S=\{s,c\}$: $(S,S')^+=\{sa,sb,cb,ct\}$ e $(S,S')^-=\{bc\}$. A capacidade é $c(s,a)+c(s,b)+c(c,b)+c(c,t)=3+2+2+3=10$. O fluxo é $f(s,a)+f(s,b)+f(c,b)+f(c,t)-f(b,c)=2+2+1+3-1=7$, igual a $f(D)$ (Teorema~\ref{teo:fluxocorte}) e menor ou igual a $10$ (Teorema~\ref{teo:cota}).

\resp{fl-maximal} Na rede residual do fluxo de valor 6, $sa$ tem $1/3$ e dá a aresta de aumento $s\to a$ com capacidade 2; $da$ tem fluxo 1 e dá a aresta de redução $a\to d$ com capacidade 1; $dt$ tem $3/5$ e dá $d\to t$ com capacidade 2. O caminho $s\to a\to d\to t$ é aumentante, com gargalo 1. O novo fluxo tem $f(s,a)=2$, $f(d,a)=0$ e $f(d,t)=4$, e valor 7: é o da Figura~\ref{fig:maximal}(b). Ele é máximo porque o corte $S=\{s,a,b,c\}$ tem capacidade $c(b,d)+c(c,t)=4+3=7$.

\resp{fl-corolario} Pelo Teorema~\ref{teo:fluxocorte}, $f(D)=f(S,S')=\sum_{(S,S')^+}f(e)-\sum_{(S,S')^-}f(e)\le\sum_{(S,S')^+}c(e)=c(S,S')$, e a igualdade vale se e somente se $f(e)=c(e)$ em $(S,S')^+$ e $f(e)=0$ em $(S,S')^-$. Se vale a igualdade, $f(D)=c(S,S')$, e pelo Teorema~\ref{teo:cota} $f$ é máximo e o corte é mínimo. Reciprocamente, se $f$ é máximo e o corte é mínimo, o Teorema~\ref{teo:maxflowmincut} dá $f(D)=c(S,S')$, e a desigualdade acima precisa ser uma igualdade.

\resp{fl-def} (a) É um fluxo: as capacidades são respeitadas, em $a$ entram 3 e saem $1+2=3$, e em $b$ entram $2+1=3$ e saem 3. O valor é $3+2=5$. (b) Não é: em $a$ entram 4 e saem só $1+2=3$, e a conservação falha. O fluxo de (a) é máximo, porque o corte $S=\{s,a\}$, com $(S,S')^+=\{sb,ab,at\}$, tem capacidade $2+1+2=5$.

\resp{fl-ff} Os caminhos aumentantes são $s,a,c,t$ (gargalo 4), $s,b,d,t$ (gargalo 5), $s,a,b,d,t$ (gargalo 1) e $s,a,c,d,t$ (gargalo 1). O fluxo máximo vale $4+5+1+1=11$. Ao final, os vértices alcançáveis a partir de $s$ na rede residual são $S=\{s,a,b\}$, e o corte mínimo tem $(S,S')^+=\{ac,bd\}$, de capacidade $5+6=11$.

\resp{fl-corte} Os cortes são $S=\{s\}$, com capacidade $3+2=5$; $S=\{s,a\}$, com $2+1+2=5$; $S=\{s,b\}$, com $3+3=6$; e $S=\{s,a,b\}$, com $2+3=5$. O fluxo máximo vale 5, e o corte mínimo não é único: há três cortes de capacidade 5.

\resp{fl-aumenta} Não: o corte $S=\{s\}$ não contém $at$ e continua com capacidade 5, e o fluxo máximo continua 5. Com capacidades inteiras, aumentar em 1 a capacidade de uma aresta $e$ aumenta o fluxo máximo (em exatamente 1) se e somente se $e$ pertence a $(S,S')^+$ de \emph{todos} os cortes mínimos: nesse caso todos os cortes mínimos ganham 1 e os demais já tinham capacidade pelo menos o mínimo mais 1; se algum corte mínimo não contém $e$, ele continua com a mesma capacidade.

\resp{fl-residualex} Arestas residuais: $s\to b$ (1) e $b\to s$ (2), de $sb$; $a\to s$ (2), de $sa$, saturada; $a\to b$ (1), de $ab$, sem fluxo; $a\to t$ (1) e $t\to a$ (2), de $at$; $t\to b$ (2), de $bt$, saturada. A partir de $s$ só se alcança $b$, e de $b$ só se volta a $s$: não há caminho aumentante, e o fluxo, de valor 4, é máximo. O corte é $S=\{s,b\}$, com $(S,S')^+=\{sa,bt\}$, capacidade $2+2=4$, e $(S,S')^-=\{ab\}$, com fluxo zero.

\resp{fl-integral} Por indução nos aumentos: o fluxo inicial (zero) é inteiro; se o fluxo é inteiro, as capacidades residuais $c-f$ e $f$ são inteiras, o gargalo é inteiro e o novo fluxo também. Como cada aumento acrescenta pelo menos 1, o algoritmo termina, e pelo Teorema~\ref{teo:aumentante} o fluxo final é máximo e inteiro. Podem existir outros fluxos máximos não inteiros, mas o algoritmo, partindo de zero, nunca os produz. Exemplo: com os caminhos $s\to a\to t$ e $s\to b\to t$ de capacidade 1 e as arestas $ab$ e $ba$ de capacidade 1, o fluxo que manda $\frac12$ por $ab$ e $\frac12$ por $ba$ (além de 1 em cada caminho) também é máximo, de valor 2.

\resp{fl-pseudo} Escolhendo alternadamente os caminhos $s,a,b,t$ e $s,b,a,t$ (este usando a aresta de redução $b\to a$), cada aumento tem gargalo 1, por causa da aresta $ab$, e são necessários $2C$ aumentos para chegar ao máximo, $2C$. O Edmonds e Karp escolhe os caminhos mais curtos, $s,a,t$ e $s,b,t$, e termina com 2 aumentos. Como $C$ é escrito com $O(\log C)$ bits, $O(m\,f^*)=O(mC)$ é exponencial no tamanho da entrada: o Ford e Fulkerson com escolha arbitrária é pseudopolinomial.

\resp{fl-ek} Seja $\delta(v)$ a distância de $s$ a $v$ na rede residual. Um aumento por um caminho mais curto remove da rede residual pelo menos uma aresta $(u,v)$ do caminho (a de gargalo) e só acrescenta arestas $(v,u)$ que voltam ao longo do caminho, de $\delta+1$ para $\delta$; arestas assim não criam atalhos, e as distâncias não diminuem. Para $(u,v)$ voltar a ser gargalo, ela precisa reaparecer, o que exige um caminho mais curto usando $(v,u)$; com $\delta(v)=\delta(u)+1$ antes e $\delta(u)=\delta(v)+1$ depois, $\delta(u)$ aumenta pelo menos 2 entre duas vezes. Como as distâncias vão até $n$, cada aresta é gargalo $O(n)$ vezes, e há $O(nm)$ aumentos, cada um em $O(m)$: total $O(nm^2)$.

\resp{fl-fluxot} Aplique o Teorema~\ref{teo:fluxocorte} ao corte $S=V\setminus\{t\}$: $(S,S')^+$ é o conjunto das arestas que chegam a $t$, e $(S,S')^-$ é vazio, porque $t$ é um sumidouro. Logo $\sum_xf(x,t)=f(S,S')=f(D)=\sum_zf(s,z)$.

\resp{fl-emp} A rede tem as arestas $sa$, $sb$, $sc$, $a1$, $a2$, $b1$, $c2$, $c3$, $1t$, $2t$ e $3t$, todas de capacidade 1. O fluxo máximo vale 3, por exemplo com fluxo 1 em $b1$, $a2$ e $c3$, o que dá o emparelhamento perfeito $\{b1,a2,c3\}$.

\resp{fl-menger} Com capacidade 1 em todas as arestas, o fluxo máximo vale 2: por exemplo, os caminhos $s,a,t$ e $s,b,d,t$ (ou $s,b,c,t$) não têm arestas em comum. O corte $S=\{s\}$ tem capacidade 2: removendo $sa$ e $sb$, $t$ fica separado de $s$. Não há 3 caminhos disjuntos, pois só duas arestas saem de $s$.

\resp{fl-vertices} Troque cada vértice $v$ por dois vértices, $v_{in}$ e $v_{out}$, ligados por uma aresta $v_{in}\to v_{out}$ de capacidade $c(v)$; toda aresta que chegava a $v$ passa a chegar a $v_{in}$, e toda aresta que saía de $v$ passa a sair de $v_{out}$. Tudo o que passa por $v$ atravessa a aresta $v_{in}\to v_{out}$, e um fluxo na rede nova corresponde a um fluxo na original que respeita as capacidades dos vértices.

\resp{fl-multi} Crie uma superorigem $s^*$ com arestas $s^*\to F_1$, $s^*\to F_2$ e $s^*\to F_3$ de capacidades 10, 15 e 20, e um superdestino $t^*$ com arestas $P_1\to t^*$ e $P_2\to t^*$ de capacidades 25 e 30; as estradas mantêm as suas capacidades. O fluxo máximo de $s^*$ a $t^*$ é o maior total escoável, e ele é no máximo $10+15+20=45$, a capacidade do corte $\{s^*\}$.

\resp{fl-nd} Troque cada ligação $\{x,y\}$ de capacidade $c$ por duas arestas, $x\to y$ e $y\to x$, ambas de capacidade $c$, e calcule o fluxo máximo. Se num fluxo máximo $f(x,y)>0$ e $f(y,x)>0$, subtraia $\min\{f(x,y),f(y,x)\}$ das duas: a conservação e o valor não mudam, e passa a haver fluxo num só sentido, que não ultrapassa $c$.

\resp{fl-konig} Com capacidade 1 em $s\to x$ e $y\to t$ e $\infty$ em $x\to y$, um corte $(S,S')$ de capacidade finita não tem arestas $x\to y$ com $x\in S$ e $y\in S'$. Seja $C=(X\setminus S)\cup(Y\cap S)$. Uma aresta $xy$ do grafo com $x\in S$ tem $y\in S$, e então $y\in C$; se $x\notin S$, então $x\in C$. Logo $C$ é uma cobertura, e $|C|=c(S,S')$: as arestas $s\to x$ com $x\in S'$ e as arestas $y\to t$ com $y\in S$. Reciprocamente, uma cobertura dá um corte desse tipo com a mesma capacidade. Pelo teorema do fluxo máximo e corte mínimo, $\nu(G)=$ fluxo máximo $=$ corte mínimo $=\tau(G)$.

\resp{fl-escala} Rede: $s\to A$ com capacidade 2 para cada aluno $A$; $A\to O$ com capacidade 1 para cada optativa $O$ que $A$ indicou; $O\to t$ com capacidade $c_O$. Pelo teorema da integralidade, o algoritmo devolve um fluxo inteiro; as arestas $A\to O$ com fluxo 1 são as matrículas. Cada aluno faz no máximo 2 (capacidade de $s\to A$), não se matricula duas vezes na mesma optativa (capacidade 1) e cada optativa recebe no máximo $c_O$ alunos. O valor do fluxo é o número de matrículas, e ele é máximo.

\resp{fl-decomp} Enquanto houver fluxo positivo saindo de $s$, siga, a partir de $s$, arestas com fluxo positivo: pela conservação sempre há como continuar, até chegar a $t$ ou repetir um vértice, formando um ciclo. Subtraia de todas as arestas do caminho (ou ciclo) o menor fluxo entre elas; isso preserva a conservação e zera pelo menos uma aresta. Quando não sair mais nada de $s$, o que sobra são ciclos, retirados do mesmo modo. Como cada parcela zera uma aresta, são no máximo $m$ parcelas. O fluxo de valor 7 se decompõe em $s,c,t$ (3), $s,b,d,t$ (2), $s,a,b,d,t$ (2) e o ciclo $b,c,b$ (1), que não contribui para o valor.
